\lecture{Лекция 1}{}

\section{Задачи линейного программирования}

\subsection{Угловые точки и многогранные множества}

\begin{definition}\label{def:extreme-point}
Пусть $S \subset \R^n$ — выпуклое множество. Точка $x \in S$ называется
\textbf{угловой} (крайней) точкой множества $S$, если её нельзя представить
в виде полусуммы двух \emph{различных} точек этого множества:
\[
  \nexists\ x_1, x_2 \in S,\ x_1 \ne x_2 : \qquad x = \tfrac{1}{2}\,(x_1 + x_2).
\]
\end{definition}

\begin{idea}
Эквивалентная формулировка: $x$ — угловая точка, если из
$x = \lambda x_1 + (1-\lambda) x_2$ с $x_1, x_2 \in S$ и $\lambda \in (0,1)$
следует $x_1 = x_2 = x$. То есть $x$ не лежит \emph{внутри} никакого отрезка,
целиком содержащегося в $S$. Форма с $\lambda = \tfrac12$ не слабее общей:
если $x$ делит отрезок в отношении $\lambda : (1-\lambda)$, то, укоротив длинную
сторону, получим отрезок множества с серединой в $x$.
\end{idea}

\inputfig{figures/fig-extreme-point}

\begin{definition}\label{def:polyhedral}
\textbf{Многогранное множество} — множество, образованное пересечением конечного
числа замкнутых полупространств и гиперплоскостей.
\end{definition}

\begin{idea}
Полупространство $\{x : \langle a, x\rangle \le \beta\}$ и гиперплоскость
$\{x : \langle a, x\rangle = \beta\}$ выпуклы и замкнуты, а пересечение выпуклых
замкнутых множеств выпукло и замкнуто. Значит, многогранное множество всегда
выпукло и замкнуто — определение~\ref{def:extreme-point} к нему применимо.
Ограниченным оно быть не обязано (например, $\{x \ge 0\}$).
\end{idea}

\subsection{Три формы задачи линейного программирования}

Задачи ЛП могут быть представлены в трёх формах.

\medskip
\noindent\textbf{(1) Каноническая форма:}
\[
  \begin{cases}
    c\T[N]\, x[N] \to \min, \\
    A[M,N]\, x[N] = b[M], \\
    x[N] \ge 0 .
  \end{cases}
\]

\noindent\textbf{(2) Симметрическая форма:}
\[
  \begin{cases}
    c\T[N]\, x[N] \to \min, \\
    A[M,N]\, x[N] \ge b[M], \\
    x[N] \ge 0 .
  \end{cases}
\]

\noindent\textbf{(3) Общая форма} — часть ограничений равенства, часть неравенства,
причём условие неотрицательности наложено не на все переменные:
\[
  \begin{cases}
    c\T[N]\, x[N] \to \min, \\
    A[M_1, N]\, x[N] = b[M_1], \\
    A[M_2, N]\, x[N] \ge b[M_2], \\
    x[N_1] \ge 0, \quad x[N \setminus N_1] \ \text{— свободные переменные},
  \end{cases}
\]
где $M = M_1 \sqcup M_2$.

Ограничения (без целевой функции) задают \textbf{допустимое множество} $S$ — это
многогранное множество в смысле определения~\ref{def:polyhedral}.

\begin{idea}
Все три формы равносильны — переход между ними делается стандартными приёмами:
\begin{itemize}[nosep,left=0pt]
  \item неравенство $\langle a, x\rangle \ge \beta$ превращается в равенство
        добавлением \emph{балансной} переменной:
        $\langle a, x\rangle - s = \beta$, $s \ge 0$;
  \item равенство разбивается на пару неравенств
        $\langle a, x\rangle \ge \beta$ и $-\langle a, x\rangle \ge -\beta$;
  \item свободная переменная заменяется разностью двух неотрицательных:
        $x_j = x_j^{+} - x_j^{-}$, $x_j^{\pm} \ge 0$;
  \item максимизация — это минимизация с $c \mapsto -c$.
\end{itemize}
Поэтому вся дальнейшая теория строится только для канонической формы.
\end{idea}

\subsection{Допустимая точка и разложение по столбцам}

Далее везде рассматривается \textbf{каноническая форма}. Примем два предположения:
\[
  \abs{M} \le \abs{N}, \qquad \rank A[M,N] = \abs{M} .
\]

\begin{idea}
Первое условие ($m \le n$) означает, что уравнений не больше, чем неизвестных, —
иначе система, как правило, переопределена. Второе (строки $A$ линейно независимы)
убирает зависимые уравнения: либо они следствия остальных и их можно вычеркнуть,
либо система несовместна. Вместе они гарантируют полный ранг по строкам, а значит
из $N$ всегда можно выбрать $\abs{M}$ линейно независимых столбцов.
\end{idea}

Пусть $x_0[N]$ — \textbf{допустимая точка}, то есть $A[M,N]\,x_0[N] = b[M]$ и
$x_0[i] \ge 0$ для всех $i \in N$. Введём обозначения для индексов положительных и
нулевых компонент:
\[
  \Nzp = \{\, i \in N \mid x_0[i] > 0 \,\}, \qquad
  \Nzz = \{\, i \in N \mid x_0[i] = 0 \,\} .
\]

Распишем систему ограничений по столбцам:
\[
  \left\{
  \begin{array}{ccccccc}
    a_{11} x_1 &+& a_{12} x_2 &+ \ \dots \ +& a_{1n} x_n &=& b_1, \\
    \vdots & & \vdots & & \vdots & & \vdots \\
    a_{m1} x_1 &+& a_{m2} x_2 &+ \ \dots \ +& a_{mn} x_n &=& b_m,
  \end{array}
  \right.
\]
где столбцы обозначены $a_1[M], a_2[M], \dots, a_n[M]$. В векторном виде это
\[
  \sum_{j \in N} x[j]\, a_j[M] = b[M] ,
\]
а для допустимой точки $x_0$ нулевые слагаемые выпадают, и остаётся
\begin{equation}\label{eq:column-decomp}
  \sum_{j \in \Nzp} x_0[j]\, a_j[M] = b[M] .
\end{equation}
